\textbf{Target modules.} Table~\ref{tab:targets} compares LoRA on all six linear maps with LoRA on $\{q,v\}$ only. At rank 1, $\{q,v\}$ (512 parameters) reaches 0.577 on descending sort and 0.056 on reversal, against 0.978 and 0.799 for all-linear (1,792 parameters). At rank 4, $\{q,v\}$ (2,048 parameters) reaches 0.971 and 1.000 against 0.997 and 0.999. The gap on descending sort at rank 4 is 0.026, with the two settings differing in parameters (2,048 against 7,168), so this ablation changes the parameter count together with the targets and does not separate the two. It does show that, at tiny rank, restricting LoRA to $\{q,v\}$ (a choice common in the LoRA literature) reaches lower accuracy than all-linear in this setting.

\begin{table}[t]\centering\caption{LoRA target modules (5 seeds), mean $\pm$ std exact match.}\label{tab:targets}\resizebox{\columnwidth}{!}{\begin{tabular}{llrcc}
\toprule
Rank & Targets & Params & sort\_desc acc & reverse acc \\
\midrule
1 & all linear & \rhval{main/lora-r-1/reverse/trainable_params/mean} & \rhval{main/lora-r-1/sort_desc/adapt_acc/mean:3} $\pm$ \rhval{main/lora-r-1/sort_desc/adapt_acc/std:3} & \rhval{main/lora-r-1/reverse/adapt_acc/mean:3} $\pm$ \rhval{main/lora-r-1/reverse/adapt_acc/std:3} \\
1 & q,v only & \rhval{abl_targets/lora-r-1-q-v-only/reverse/trainable_params/mean} & \rhval{abl_targets/lora-r-1-q-v-only/sort_desc/adapt_acc/mean:3} $\pm$ \rhval{abl_targets/lora-r-1-q-v-only/sort_desc/adapt_acc/std:3} & \rhval{abl_targets/lora-r-1-q-v-only/reverse/adapt_acc/mean:3} $\pm$ \rhval{abl_targets/lora-r-1-q-v-only/reverse/adapt_acc/std:3} \\
4 & all linear & \rhval{main/lora-r-4/reverse/trainable_params/mean} & \rhval{main/lora-r-4/sort_desc/adapt_acc/mean:3} $\pm$ \rhval{main/lora-r-4/sort_desc/adapt_acc/std:3} & \rhval{main/lora-r-4/reverse/adapt_acc/mean:3} $\pm$ \rhval{main/lora-r-4/reverse/adapt_acc/std:3} \\
4 & q,v only & \rhval{abl_targets/lora-r-4-q-v-only/reverse/trainable_params/mean} & \rhval{abl_targets/lora-r-4-q-v-only/sort_desc/adapt_acc/mean:3} $\pm$ \rhval{abl_targets/lora-r-4-q-v-only/sort_desc/adapt_acc/std:3} & \rhval{abl_targets/lora-r-4-q-v-only/reverse/adapt_acc/mean:3} $\pm$ \rhval{abl_targets/lora-r-4-q-v-only/reverse/adapt_acc/std:3} \\
\bottomrule
\end{tabular}
}\end{table}

\textbf{Adaptation set size.} Table~\ref{tab:ntrain} (descending sort, seeds 0--2, same learning rates) examines low-data behaviour. With $N=300$ the means are 0.765 (full fine-tuning), 0.511 (scratch), 0.838, 0.938 and 0.621 (LoRA $r=1,4,8$); with $N=100$ they are 0.204, 0.116, 0.385, 0.337 and 0.000. Pretrained full fine-tuning is ahead of scratch at both sizes, so pretraining is useful when data are scarce, in this three-seed sample. LoRA at ranks 1 and 4 has higher means than full fine-tuning at $N=300$ and $N=100$, but we ran no test that supports an ordering: standard deviations at $N=100$ (0.15--0.18) are as large as the gaps, paired $t$-tests over the three seeds (computed from the registry rows, not reported in a table) give $p>0.05$ for all four comparisons of LoRA $r=1,4$ with full fine-tuning, and the full fine-tuning rate ($10^{-2}$) was tuned for $N=2000$ and may be too high for small $N$, so the comparison may favour LoRA unfairly. LoRA $r=4$ at $N=300$ (0.938, std 0.007) has the largest gap in the means but is also not significant. Rank 8 is at or below 0.001 in all three seeds at $N=100$ and in one of three seeds at $N=300$ (std 0.535). In these failed runs the final training loss stayed high, so they look like optimisation failures at the shared high learning rate and not memorisation; we did not test whether a lower rate repairs them. Learning rates were not retuned per $N$.

\begin{table}[t]\centering\caption{Descending-sort exact match against adaptation set size $N$ (3 seeds; $N=2000$ column from the main runs, seeds 0--2).}\label{tab:ntrain}\resizebox{\columnwidth}{!}{\begin{tabular}{lccc}
\toprule
System & $N$=100 & $N$=300 & $N$=2000 \\
\midrule
Full fine-tuning & 0.204 $\pm$ 0.027 & 0.765 $\pm$ 0.087 & 0.998 $\pm$ 0.001 \\
From scratch & 0.116 $\pm$ 0.010 & 0.511 $\pm$ 0.059 & 0.997 $\pm$ 0.000 \\
Last block & 0.075 $\pm$ 0.018 & 0.543 $\pm$ 0.069 & 0.965 $\pm$ 0.010 \\
LoRA r=1 & 0.385 $\pm$ 0.184 & 0.838 $\pm$ 0.057 & 0.978 $\pm$ 0.009 \\
LoRA r=4 & 0.337 $\pm$ 0.146 & 0.938 $\pm$ 0.007 & 0.996 $\pm$ 0.001 \\
LoRA r=8 & 0.000 $\pm$ 0.000 & 0.621 $\pm$ 0.535 & 0.998 $\pm$ 0.000 \\
\bottomrule
\end{tabular}
}\end{table}

\textbf{Learning rate.} Table~\ref{tab:lr} and Fig.~\ref{fig:lr} sweep the learning rate for full fine-tuning and LoRA $r=4$ (seeds 0--2). Full fine-tuning is stable over the range: 0.956 to 0.998 on descending sort, and at least 0.996 on reversal. LoRA is much more sensitive at the low end: at $3\times10^{-4}$ it reaches 0.517 and 0.247 (std 0.272 on reversal), and on descending sort it reaches 0.992 at $10^{-2}$ and 0.996 at its default $3\times10^{-2}$, against 0.998 for full fine-tuning at its default. Pretask retention is about 0.001 at every learning rate in the table, for both methods, including where LoRA has barely adapted.

\begin{table}[t]\centering\caption{Learning-rate sweep, seeds 0--2 (the largest rate of each system is the main-group default). ``Reverse retention'': sort-ascending exact match after adapting to reversal.}\label{tab:lr}\resizebox{\columnwidth}{!}{\begin{tabular}{llccc}
\toprule
System & LR & sort\_desc & reverse & reverse retention \\
\midrule
Full fine-tuning & 0.0003 & \rhval{sweep_lr/full-fine-tuning@lr=0.0003/sort_desc/adapt_acc/mean:3} $\pm$ \rhval{sweep_lr/full-fine-tuning@lr=0.0003/sort_desc/adapt_acc/std:3} & \rhval{sweep_lr/full-fine-tuning@lr=0.0003/reverse/adapt_acc/mean:3} $\pm$ \rhval{sweep_lr/full-fine-tuning@lr=0.0003/reverse/adapt_acc/std:3} & \rhval{sweep_lr/full-fine-tuning@lr=0.0003/reverse/pretask_acc/mean:3} $\pm$ \rhval{sweep_lr/full-fine-tuning@lr=0.0003/reverse/pretask_acc/std:3} \\
Full fine-tuning & 0.001 & \rhval{sweep_lr/full-fine-tuning@lr=0.001/sort_desc/adapt_acc/mean:3} $\pm$ \rhval{sweep_lr/full-fine-tuning@lr=0.001/sort_desc/adapt_acc/std:3} & \rhval{sweep_lr/full-fine-tuning@lr=0.001/reverse/adapt_acc/mean:3} $\pm$ \rhval{sweep_lr/full-fine-tuning@lr=0.001/reverse/adapt_acc/std:3} & \rhval{sweep_lr/full-fine-tuning@lr=0.001/reverse/pretask_acc/mean:3} $\pm$ \rhval{sweep_lr/full-fine-tuning@lr=0.001/reverse/pretask_acc/std:3} \\
Full fine-tuning & 0.003 & \rhval{sweep_lr/full-fine-tuning@lr=0.003/sort_desc/adapt_acc/mean:3} $\pm$ \rhval{sweep_lr/full-fine-tuning@lr=0.003/sort_desc/adapt_acc/std:3} & \rhval{sweep_lr/full-fine-tuning@lr=0.003/reverse/adapt_acc/mean:3} $\pm$ \rhval{sweep_lr/full-fine-tuning@lr=0.003/reverse/adapt_acc/std:3} & \rhval{sweep_lr/full-fine-tuning@lr=0.003/reverse/pretask_acc/mean:3} $\pm$ \rhval{sweep_lr/full-fine-tuning@lr=0.003/reverse/pretask_acc/std:3} \\
LoRA r=4 & 0.0003 & \rhval{sweep_lr/lora-r-4@lr=0.0003/sort_desc/adapt_acc/mean:3} $\pm$ \rhval{sweep_lr/lora-r-4@lr=0.0003/sort_desc/adapt_acc/std:3} & \rhval{sweep_lr/lora-r-4@lr=0.0003/reverse/adapt_acc/mean:3} $\pm$ \rhval{sweep_lr/lora-r-4@lr=0.0003/reverse/adapt_acc/std:3} & \rhval{sweep_lr/lora-r-4@lr=0.0003/reverse/pretask_acc/mean:3} $\pm$ \rhval{sweep_lr/lora-r-4@lr=0.0003/reverse/pretask_acc/std:3} \\
LoRA r=4 & 0.001 & \rhval{sweep_lr/lora-r-4@lr=0.001/sort_desc/adapt_acc/mean:3} $\pm$ \rhval{sweep_lr/lora-r-4@lr=0.001/sort_desc/adapt_acc/std:3} & \rhval{sweep_lr/lora-r-4@lr=0.001/reverse/adapt_acc/mean:3} $\pm$ \rhval{sweep_lr/lora-r-4@lr=0.001/reverse/adapt_acc/std:3} & \rhval{sweep_lr/lora-r-4@lr=0.001/reverse/pretask_acc/mean:3} $\pm$ \rhval{sweep_lr/lora-r-4@lr=0.001/reverse/pretask_acc/std:3} \\
LoRA r=4 & 0.003 & \rhval{sweep_lr/lora-r-4@lr=0.003/sort_desc/adapt_acc/mean:3} $\pm$ \rhval{sweep_lr/lora-r-4@lr=0.003/sort_desc/adapt_acc/std:3} & \rhval{sweep_lr/lora-r-4@lr=0.003/reverse/adapt_acc/mean:3} $\pm$ \rhval{sweep_lr/lora-r-4@lr=0.003/reverse/adapt_acc/std:3} & \rhval{sweep_lr/lora-r-4@lr=0.003/reverse/pretask_acc/mean:3} $\pm$ \rhval{sweep_lr/lora-r-4@lr=0.003/reverse/pretask_acc/std:3} \\
LoRA r=4 & 0.01 & \rhval{sweep_lr/lora-r-4@lr=0.01/sort_desc/adapt_acc/mean:3} $\pm$ \rhval{sweep_lr/lora-r-4@lr=0.01/sort_desc/adapt_acc/std:3} & \rhval{sweep_lr/lora-r-4@lr=0.01/reverse/adapt_acc/mean:3} $\pm$ \rhval{sweep_lr/lora-r-4@lr=0.01/reverse/adapt_acc/std:3} & \rhval{sweep_lr/lora-r-4@lr=0.01/reverse/pretask_acc/mean:3} $\pm$ \rhval{sweep_lr/lora-r-4@lr=0.01/reverse/pretask_acc/std:3} \\
\bottomrule
\end{tabular}
}\end{table}

\begin{figure}[t]\centering\includegraphics[width=\columnwidth]{lr_sweep.pdf}\caption{Exact match against learning rate (mean $\pm$ std over 3 seeds).}\label{fig:lr}\end{figure}

\textbf{Adaptation speed.} Fig.~\ref{fig:curves} (descending sort, seeds 0--2) shows that full fine-tuning reaches a mean target accuracy of 0.9 first (step 350), then last block (500) and LoRA at ranks 1 and 4 (550 each); the retention panel is zero for all systems at every evaluation point from step 50 on. Per-step differences are within the seed band shown, so we do not rank the three slower systems on speed.

\begin{figure}[t]\centering\includegraphics[width=\columnwidth]{curves_custom.pdf}\caption{Adaptation to descending sort (left) and retention of ascending sort (right) during training; mean with min--max band over seeds 0--2.}\label{fig:curves}\end{figure}
